% Generated draft: requires mathematical and source review.
\documentclass[11pt]{article}
\usepackage[margin=25mm]{geometry}
\usepackage{fontspec}
\setmainfont{Latin Modern Roman}
\usepackage{amsmath,amssymb,mathtools,unicode-math}
\setmathfont{Latin Modern Math}
\usepackage{enumitem,needspace,xurl,hyperref}
\setlength{\parindent}{0pt}
\setlength{\parskip}{5pt}
\setlength{\emergencystretch}{3em}
\newcommand{\N}{\mathbb{N}}
\newcommand{\Z}{\mathbb{Z}}
\newcommand{\Q}{\mathbb{Q}}
\newcommand{\R}{\mathbb{R}}
\newcommand{\C}{\mathbb{C}}
\newcommand{\F}{\mathbb{F}}
\newcommand{\A}{\mathbb{A}}
\newcommand{\PP}{\mathbb{P}}
\newcommand{\E}{\mathbb{E}}
\newcommand{\eps}{\varepsilon}
\newcommand{\dd}{\,\mathrm{d}}
\begin{document}
\section*{UnsolvedMath ID: 701231}
\textbf{Source problem code:} catalogue entry
\subsection{Definitions and mathematical statement}
Let \(k\) be an algebraically closed field of characteristic \(p>0\), and let \(G\) be a semisimple algebraic group over \(k\). Set \(\mathfrak g=\operatorname{Lie}(G)\), equipped with its restricted Lie algebra operation \(x\mapsto x^{[p]}\), and let
\[
u=u(\mathfrak g)=U(\mathfrak g)\big/\langle x^p-x^{[p]}:x\in\mathfrak g\rangle
\]
be its restricted enveloping algebra. A left \(u\)-module is projective if it is a direct summand of a free left \(u\)-module, and is indecomposable if it is nonzero and cannot be expressed as a direct sum of two nonzero submodules. A rational representation of \(G\) differentiates to a restricted representation of \(\mathfrak g\), hence gives a \(u\)-module, denoted \(V|_u\).

\textbf{Question (J. E. Humphreys and D. N. Verma; proposed in the notebook by S. Donkin).} For every such \(k\) and \(G\), and every projective indecomposable left \(u\)-module \(P\), does there exist a rational \(G\)-module \(V\) such that
\[
V|_u\cong P
\]
as \(u\)-modules?

The supplied source attributes to Curtis the statement that every irreducible restricted \(u\)-module is obtained by differentiating a rational \(G\)-module. It further attributes to Jantzen an affirmative answer to the question when \(p\geq 2h-2\), where it calls \(h\) the Coxeter number of \(G\). These are supplied background assertions, not independently verified here. The source does not specify how \(h\) is chosen when the root system of \(G\) has components with different Coxeter numbers. It also supplies no simply connected hypothesis; the question above retains its stated semisimple scope, whose compatibility with the attributed background results requires review.
\subsection{Short English statement}
Does every projective indecomposable module for the restricted enveloping algebra of a semisimple algebraic group in positive characteristic extend to a rational representation of that same group? The source reports an affirmative answer for \(p\geq 2h-2\), subject to clarification of \(h\) and the precise hypotheses of the cited results.
\subsection{Sources}
\begin{itemize}
\item Local UnsolvedMath registry, permanent ID \(701231\), source code \texttt{catalogue entry}, title \emph{Kourovka Notebook Problem 16.39}. The supplied record labels the problem \texttt{open}; this is a registry classification, not an independently verified assessment.
\item \emph{The Kourovka Notebook (Unsolved Problems in Group Theory)}, problem 16.39, 16th issue (2006); also identified in the supplied background with the 21st issue edition (2026), \texttt{arXiv:1401.0300v48}. The collection description names E. I. Khukhro and V. D. Mazurov as editors.
\item The registry identifies S. Donkin as proposer, with proposed year 2006. The statement itself attributes the question to J. E. Humphreys and D. N. Verma.
\item The supplied statement attributes the irreducible-module extension result to Curtis and the range \(p\geq 2h-2\) to results of Jantzen. No specific publications or theorem references for these assertions are supplied.
\end{itemize}
\end{document}
